\begin{figure}[htbp]
\centering
\begin{tikzpicture}

% ========== Main Panel: Eigenvector Visualization ==========
\begin{scope}[xshift=0cm]
    \node[font=\bfseries\large] at (4, 5.5) {Eigenvectors: Directions Preserved Under Transformation};

    % Axes
    \draw[->, thick] (-3.5, 0) -- (5.5, 0) node[right] {$x$};
    \draw[->, thick] (0, -3.5) -- (0, 5) node[above] {$y$};

    % Grid
    \draw[help lines, UofSC50Black!30] (-3, -3) grid (5, 4.5);

    % Matrix A = [2 1; 1 2] has eigenvalues 3 and 1
    % Eigenvector for lambda=3: v1 = [1; 1] (direction 45 degrees)
    % Eigenvector for lambda=1: v2 = [1; -1] (direction -45 degrees)

    % ========== Eigenvector 1: lambda = 3 ==========
    % Original eigenvector v1 (shorter, dashed)
    \draw[->, line width=2pt, UofSCGarnet, dashed] (0, 0) -- (1.2, 1.2);
    \node[UofSCGarnet, font=\bfseries, above left] at (0.9, 0.9) {$\vecv_1$};

    % Transformed Av1 = 3*v1 (longer, solid)
    \draw[->, line width=2.5pt, UofSCGarnet] (0, 0) -- (3.6, 3.6);
    \node[UofSCGarnet, font=\bfseries, above right] at (3.4, 3.4) {$\mA\vecv_1 = 3\vecv_1$};

    % Scaling annotation for v1
    \draw[<->, UofSC70Black, thick] (1.4, 1.0) -- (3.4, 3.0);
    \node[UofSC70Black, font=\footnotesize, fill=white, inner sep=1pt] at (2.4, 2.0) {$\times 3$};

    % ========== Eigenvector 2: lambda = 1 ==========
    % Original eigenvector v2 (dashed)
    \draw[->, line width=2pt, UofSCAtlantic, dashed] (0, 0) -- (1.5, -1.5);
    \node[UofSCAtlantic, font=\bfseries, below right] at (1.2, -1.2) {$\vecv_2$};

    % Transformed Av2 = 1*v2 (same length, solid - overlaps)
    \draw[->, line width=2.5pt, UofSCAtlantic] (0, 0) -- (1.5, -1.5);
    \node[UofSCAtlantic, font=\bfseries, below left] at (0.6, -1.8) {$\mA\vecv_2 = 1 \cdot \vecv_2$};

    % Scaling annotation for v2
    \node[UofSC70Black, font=\footnotesize, fill=white, inner sep=2pt] at (2.5, -1.5) {$\times 1$ (unchanged)};

    % ========== Non-eigenvector for comparison ==========
    % Original vector w (not an eigenvector)
    \draw[->, line width=1.5pt, UofSCHorseshoe, dashed] (0, 0) -- (2, 0.5);
    \node[UofSCHorseshoe, font=\bfseries, above] at (1.8, 0.5) {$\vecw$};

    % Transformed Aw (direction changes!)
    % A*[2; 0.5] = [2*2 + 1*0.5; 1*2 + 2*0.5] = [4.5; 3]
    \draw[->, line width=2pt, UofSCHorseshoe] (0, 0) -- (4.5, 3);
    \node[UofSCHorseshoe, font=\bfseries, above left] at (4.3, 2.8) {$\mA\vecw$};

    % Direction change annotation
    \draw[->, UofSCRose, thick, dashed] (2.2, 0.7) to[bend left=20] (3.8, 2.2);
    \node[UofSCRose, font=\footnotesize, fill=white, inner sep=1pt] at (3.5, 1.2) {direction changes!};

    % Origin
    \node[below left, font=\small] at (0, 0) {$O$};
    \node[circle, fill=UofSCBlack, minimum size=5pt, inner sep=0pt] at (0, 0) {};

    % ========== Matrix and Eigenvalue Info Box ==========
    \node[draw=UofSC70Black, thick, fill=white, align=left, font=\footnotesize,
          text width=5.5cm, inner sep=8pt] at (8, 2) {
        \textbf{Matrix:}\\[2pt]
        $\mA = \begin{bmatrix} 2 & 1 \\ 1 & 2 \end{bmatrix}$\\[8pt]
        \textbf{Characteristic equation:}\\[2pt]
        $\det(\mA - \lambda\mI) = \lambda^2 - 4\lambda + 3 = 0$\\[8pt]
        \textbf{Eigenvalues:}\\[2pt]
        $\lambda_1 = 3$, \quad $\lambda_2 = 1$\\[8pt]
        \textbf{Eigenvectors:}\\[2pt]
        $\vecv_1 = \begin{bmatrix} 1 \\ 1 \end{bmatrix}$ for $\lambda_1 = 3$\\[4pt]
        $\vecv_2 = \begin{bmatrix} 1 \\ -1 \end{bmatrix}$ for $\lambda_2 = 1$
    };

    % ========== Key Concept Box ==========
    \node[draw=UofSCGarnet, thick, fill=UofSCGarnet!8, align=left, font=\footnotesize,
          text width=5.5cm, inner sep=8pt] at (8, -2) {
        \textbf{Key Insight:}\\[4pt]
        Eigenvectors are special directions that remain\\
        \textit{parallel to themselves} under transformation.\\[4pt]
        The eigenvalue $\lambda$ is the \textbf{scaling factor}:\\[2pt]
        \textbullet\ $|\lambda| > 1$: stretching\\
        \textbullet\ $|\lambda| < 1$: compression\\
        \textbullet\ $\lambda < 0$: reversal\\[4pt]
        Non-eigenvectors (like $\vecw$) change direction.
    };

\end{scope}

\end{tikzpicture}
\caption{Geometric visualization of eigenvalues and eigenvectors. For matrix $\mA$, eigenvectors $\vecv_1$ and $\vecv_2$ (dashed) are special directions that remain parallel to themselves after transformation (solid). The eigenvalue is the scaling factor: $\vecv_1$ is scaled by $\lambda_1 = 3$ (stretched), while $\vecv_2$ is scaled by $\lambda_2 = 1$ (unchanged). In contrast, a non-eigenvector $\vecw$ both rotates and scales under transformation---its direction is not preserved. This geometric property is fundamental to understanding system stability: eigenvalues determine how state components grow or decay along eigenvector directions.}
\label{fig:eigenvalue_visualization}
\end{figure}
